\chapter{Gaussian and multinomial inversion}
\label{chap:gaussian-multinomial-inversion}

\section{Polynomial structure and the positive real inverse}
\label{sec:gaussian-polynomial}

Fix integers \(0<k<n\), put \(b=n-k\), \(d=kb\), and write
\[
  G_{n,k}(q):=\GaussianBinomial{n}{k}{q}.
\]
Although the quotient in \cref{sec:normalizations} is useful away from roots
of unity, the polynomial itself is the primary object.

\begin{proposition}[Positive palindromic Gaussian polynomial]
\label{prop:gq-positive-palindromic}
The function \(G_{n,k}\) is a monic polynomial of degree \(d\), with
nonnegative integer coefficients, constant coefficient one, and
\[
  G_{n,k}(q)=q^dG_{n,k}(q^{-1})\qquad(q\ne0).
\]
Its coefficient of \(q\) is one.
\end{proposition}

\begin{proof}
The finite \(q\)-binomial recurrence is
\[
  G_{n,k}(q)
  =G_{n-1,k}(q)+q^{n-k}G_{n-1,k-1}(q).
\]
It follows directly from the product quotient: after putting the two terms
over the denominator
\(\QPochhammer{q}{q}{k}\QPochhammer{q}{q}{n-k}\), the numerator becomes
\[
  \QPochhammer{q}{q}{n-1}\bigl((1-q^{n-k})+q^{n-k}(1-q^k)\bigr)
  =\QPochhammer{q}{q}{n-1}(1-q^n).
\]
Starting from \(G_{m,0}=G_{m,m}=1\), the recurrence proves inductively that
all coefficients are nonnegative integers.  The same induction gives degree
\(k(n-k)\), monicity, and constant coefficient one.

For reciprocity, replace \(q\) by \(q^{-1}\) in
\[
  G_{n,k}(q)=\prod_{j=1}^{k}\frac{1-q^{b+j}}{1-q^j}.
\]
Since \(1-q^{-m}=-q^{-m}(1-q^m)\), the signs cancel and the accumulated
power is
\[
  -\sum_{j=1}^{k}(b+j)+\sum_{j=1}^{k}j=-kb=-d.
\]
Thus \(G_{n,k}(q^{-1})=q^{-d}G_{n,k}(q)\).
Finally, the recurrence or the partition interpretation shows that the
coefficient of \(q\) is one: there is exactly one partition of \(1\) in a
nonempty \(k\times b\) rectangle.
\end{proof}

\begin{remark}[Exact Lean crosswalk]
The universal-polynomial theorem
\texttt{gaussianBinomial\_\allowbreak{}eq\_\allowbreak{}eval}%
\textsubscript{2}\allowbreak{}\texttt{\_universal} supplies the nonnegative
integral coefficients.  The direct structure theorems of the
zero-definition, fourteen-theorem
\nolinkurl{GaussianBinomialPalindromic} module---
\nolinkurl{gaussianBinomial_natDegree},
\nolinkurl{gaussianBinomial_monic},
\nolinkurl{coeff_gaussianBinomial_zero},
\nolinkurl{reflect_gaussianBinomial}, and
\nolinkurl{coeff_gaussianBinomial_reflect}---give exact degree,
monicity, constant coefficient one, reflection, and coefficient
palindromicity for every admissible column over generic commutative
semirings.  The same API gives the top coefficient and the division-free
mean identity \nolinkurl{two_mul_derivative_gaussianBinomial_eval_one}.
Its new declarations
\nolinkurl{coeff_gaussianBinomial_one_of_pos_of_lt} and
\nolinkurl{coeff_gaussianBinomial_one} prove, over every commutative
semiring, that the coefficient of \(q\) is one for \(0<k<n\), and classify
it totally as zero on the lower edge, on the diagonal, and above the
diagonal.
The companion zero-definition, five-theorem
\nolinkurl{GaussianBinomialPolynomialStructure} module exports the same
degree, monic, constant, reflection, and palindrome facts specialized to
\(\NaturalNumbers[X]\).  Together these declarations give an exact Lean
counterpart of \cref{prop:gq-positive-palindromic}.
\end{remark}

\begin{theorem}[Global positive-base inverse]
\label{thm:gq-positive-inverse}
The restriction of \(G_{n,k}\) to \([0,\infty)\) is a continuous strictly
increasing bijection onto \([1,\infty)\).  Consequently, for every real
\(y>1\) there is exactly one \(q>0\) satisfying \(G_{n,k}(q)=y\).
\end{theorem}

\begin{proof}
Write \(G_{n,k}(q)=1+\sum_{r=1}^{d}p_rq^r\).  By
\cref{prop:gq-positive-palindromic}, \(p_r\ge0\) and \(p_1=1\).  Hence
\[
  G_{n,k}'(q)=\sum_{r=1}^{d}rp_rq^{r-1}\ge1
  \qquad(q\ge0).
\]
The function is therefore strictly increasing, starts at \(G_{n,k}(0)=1\),
and tends to infinity because it is monic of positive degree.  Continuity and
the intermediate value theorem give existence; strict increase gives
uniqueness.
\end{proof}

\section{\texorpdfstring{The inverse at \(q=0\)}{The inverse at q=0}}
\label{sec:gaussian-zero}

Let \(p_{k,b}(r)\) denote the number of partitions of \(r\) whose Ferrers
diagram fits in a \(k\times b\) rectangle.  The recurrence in the preceding
proof gives the standard generating identity
\[
  G_{n,k}(q)=\sum_{r=0}^{d}p_{k,b}(r)q^r.
\]
In particular, \(p_{k,b}(0)=p_{k,b}(1)=1\).  Put
\[
  c_2=\IndicatorOf{k\ge2}+\IndicatorOf{b\ge2},
  \qquad
  c_3=\IndicatorOf{b\ge3}
      +\IndicatorOf{k\ge2,\,b\ge2}
      +\IndicatorOf{k\ge3}.
\]
These are exactly \(p_{k,b}(2)\) and \(p_{k,b}(3)\).

\begin{theorem}[Gaussian reversion at zero]
\label{thm:gq-zero-inverse}
Let \(u=y-1\).  The positive inverse from
\cref{thm:gq-positive-inverse} has the convergent local expansion
\[
  q=u-c_2u^2+(2c_2^2-c_3)u^3+\BigO(u^4).
\]
More generally,
\[
  [u^m]q(u)
  =\frac1m[z^{m-1}]
    \left(\frac{z}{G_{n,k}(z)-1}\right)^m
  \qquad(m\ge1).
\]
\end{theorem}

\begin{proof}
Set \(F(z)=G_{n,k}(z)-1\).  Then \(F(0)=0\) and \(F'(0)=1\), so the
analytic inverse function theorem gives a convergent inverse near zero.
The Lagrange--Bürmann formula from
\cref{thm:uq-analytic-lagrange} gives the coefficient identity.  For the
first terms, substitute
\[
  u=q+c_2q^2+c_3q^3+\BigO(q^4),
  \qquad
  q=u+A u^2+B u^3+\BigO(u^4).
\]
Comparison of \(u^2\) and \(u^3\) gives \(A=-c_2\) and
\(B=2c_2^2-c_3\).
\end{proof}

\section{\texorpdfstring{Complete expansion and inversion at \(q=1\)}{Complete expansion and inversion at q=1}}
\label{sec:gaussian-one}

Let
\[
  S_m(N):=\sum_{j=1}^{N}j^m,
  \qquad
  \Delta_m(n,k):=S_m(n)-S_m(k)-S_m(n-k).
\]

\begin{lemma}[The elementary Bernoulli factor]
\label{lem:gq-bernoulli-factor}
On \(|x|<2\pi\), let \(\ell(x)\) be the unique holomorphic logarithm of
\((1-\EulerE^{-x})/x\) normalized by \(\ell(0)=0\).  Then
\[
  \ell(x)
  =-\frac{x}{2}
   +\sum_{r\ge1}\frac{B_{2r}}{2r(2r)!}x^{2r}.
\]
\end{lemma}

\begin{proof}
The generating function
\[
  \frac{x}{\EulerE^x-1}
  =\sum_{m\ge0}B_m\frac{x^m}{m!}
\]
has nearest nonzero singularities at \(\pm2\pi\ImaginaryUnit\).  Therefore, for
\(|x|<2\pi\),
\[
  \frac1{\EulerE^x-1}-\frac1x
  =-\frac12+\sum_{r\ge1}B_{2r}\frac{x^{2r-1}}{(2r)!}.
\]
The quotient \((1-\EulerE^{-x})/x\), extended by the value one at zero, is
zero-free in the open disk because its first nonzero zeros are
\(\pm2\pi\ImaginaryUnit\).  Hence its normalized holomorphic logarithm
exists there.  The left side above is its derivative.  Termwise integration
from \(0\) to \(x\) gives the formula, with zero integration constant by the
normalization.
\end{proof}

\begin{theorem}[Complete Gaussian logarithmic expansion]
\label{thm:gq-one-log}
Let \(q=\EulerE^{-t}\).  For \(|t|<2\pi/n\), take the holomorphic logarithm
of the normalized quotient whose value at \(t=0\) is zero.  Then
\[
  \log\frac{G_{n,k}(\EulerE^{-t})}{\binom nk}
  =-\frac d2t
   +\sum_{r\ge1}
    \frac{B_{2r}\Delta_{2r}(n,k)}
         {2r(2r)!}\,t^{2r}.
\]
If
\[
  \lambda_1=-\frac d2,\qquad
  \lambda_{2r}=\frac{B_{2r}}{2r}\Delta_{2r}(n,k),
  \qquad
  \lambda_{2r+1}=0\quad(r\ge1),
\]
and the complete exponential Bell polynomials \(\mathcal B_m\) are defined by
\[
  \exp\left(\sum_{j\ge1}x_j\frac{z^j}{j!}\right)
  =\sum_{m\ge0}\mathcal B_m(x_1,\ldots,x_m)\frac{z^m}{m!},
\]
then
\[
  G_{n,k}(\EulerE^{-t})
  =\binom nk\sum_{m\ge0}
    \mathcal B_m(\lambda_1,\ldots,\lambda_m)\frac{t^m}{m!}.
\]
\end{theorem}

\begin{proof}
Use
\[
  G_{n,k}(\EulerE^{-t})
  =\prod_{j=1}^{k}
    \frac{1-\EulerE^{-(b+j)t}}{1-\EulerE^{-jt}}.
\]
Apply \cref{lem:gq-bernoulli-factor} to every numerator and denominator.
The logarithms of the scale factors \((b+j)/j\) sum to
\(\log\binom nk\).  The linear difference is
\[
  \sum_{j=1}^{k}(b+j-j)=kb=d,
\]
and the even differences are precisely
\(\Delta_{2r}(n,k)\).  Every factor expansion converges when
\(|(b+j)t|<2\pi\); the sufficient common disk is \(|t|<2\pi/n\).
Exponentiating the resulting series and invoking the defining generating
identity for \(\mathcal B_m\) proves the second formula.
\end{proof}

\begin{corollary}[Literal Taylor coordinate \(h=q-1\)]
\label{cor:gq-one-h}
As \(h\to0\),
\[
  \frac{G_{n,k}(1+h)}{\binom nk}
  =1+\frac d2h
   +\frac{d(3d+n-5)}{24}h^2
   +\frac{d(d-2)(d+n-3)}{48}h^3
   +\BigO(h^4).
\]
\end{corollary}

\begin{proof}
In \cref{thm:gq-one-log}, substitute
\[
  t=-\log(1+h)=-h+\frac{h^2}{2}-\frac{h^3}{3}
  +\BigO(h^4),
\]
use \(\Delta_2=d(n+1)\), and exponentiate through degree three.
All discarded products have degree at least four.  Straight collection gives
the displayed coefficients.
\end{proof}

Put \(v=\log(y/\binom nk)\).  The linear term in
\cref{thm:gq-one-log} is nonzero.

\begin{corollary}[Ordinary Gaussian inverse at \(q=1\)]
\label{cor:gq-one-inverse}
On the local branch through \(q=1\),
\[
  t=-\frac{2v}{d}
    +\frac{\Delta_2}{3d^3}v^2
    -\frac{\Delta_2^2}{9d^5}v^3
    +\BigO(v^4),
  \qquad q=\EulerE^{-t}.
\]
The coefficient of \(v^4\) is the first one involving \(\Delta_4\).
\end{corollary}

\begin{proof}
Write
\[
  v=c_1t+c_2t^2+c_4t^4+\cdots,
  \qquad c_1=-d/2,\quad c_2=\Delta_2/24.
\]
There is no cubic term.  Substituting
\(t=Av+Bv^2+Cv^3+\cdots\) and comparing coefficients gives
\[
  A=c_1^{-1},\quad
  B=-c_2c_1^{-3},\quad
  C=2c_2^2c_1^{-5},
\]
which are the displayed values.  Since \(c_4\) first contributes in degree
four, \(\Delta_4\) cannot occur earlier.
\end{proof}

\section{Centered palindromic inversion}
\label{sec:centered-palindromic}

The ordinary inverse above retains the deterministic reciprocity drift.  Once
that drift is removed, \(q=1\) becomes a fold.

\begin{theorem}[Centered palindromic inverse]
\label{thm:gq-centered-palindromic}
Let
\[
  P(q)=\sum_{j=0}^{D}c_jq^j
\]
be nonconstant, with \(c_j\ge0\), \(c_j=c_{D-j}\), and
\(c_0=c_D>0\).  Define
\[
  H_P(t):=\EulerE^{Dt/2}\frac{P(\EulerE^{-t})}{P(1)}.
\]
Then \(H_P\) is even, \(H_P(0)=1\), strictly increases on
\((0,\infty)\), and tends to infinity there.  For every \(h>1\), the equation
\(H_P(t)=h\) therefore has exactly two real solutions
\(t=\pm E_P(h)\).  At \(h=1\) the two branches meet in a square-root fold.
\end{theorem}

\begin{proof}
Let \(X\) take value \(j\) with probability \(c_j/P(1)\), and set
\(Z=X-D/2\).  Palindromy says \(Z\) and \(-Z\) have the same law.  Hence
\[
  H_P(t)=\Expectation(\EulerE^{-tZ})=\Expectation(\cosh(tZ)),
\]
which is even and equals one at zero.  For \(t>0\),
\[
  H_P'(t)=\Expectation\bigl(Z\sinh(tZ)\bigr)>0.
\]
Indeed, the integrand is nonnegative and is positive whenever \(Z\ne0\);
the endpoint masses \(c_0=c_D>0\) ensure that this happens with positive
probability.  Those same endpoints yield
\[
  H_P(t)\ge\frac{c_0}{P(1)}\EulerE^{Dt/2}\longrightarrow\infty.
\]
This proves the global branch count.

Finally,
\[
  H_P(t)=1+\frac{\Expectation(Z^2)}2t^2+\BigO(t^4),
\]
and \(\Expectation(Z^2)>0\) by nondegeneracy.  Thus
\[
  t=\pm\sqrt{\frac{2(h-1)}{\Expectation(Z^2)}}\,
  \bigl(1+\BigO(h-1)\bigr),
\]
which is exactly a square-root fold.
\end{proof}

Let \(\kappa_{2r}\) be the even cumulants of \(Z\), so that
\[
  u:=\log H_P(t)
  =\frac{\kappa_2}{2!}t^2
   +\frac{\kappa_4}{4!}t^4
   +\frac{\kappa_6}{6!}t^6+\cdots.
\]

\begin{proposition}[Centered inverse through \(u^{5/2}\)]
\label{prop:gq-centered-puiseux}
As \(u\downarrow0\),
\[
  t=\pm\sqrt{\frac{2u}{\kappa_2}}
  \left[
    1-\frac{\kappa_4}{12\kappa_2^2}u
    +\left(
       \frac{7\kappa_4^2}{288\kappa_2^4}
       -\frac{\kappa_6}{180\kappa_2^3}
     \right)u^2
    +\BigO(u^3)
  \right].
\]
\end{proposition}

\begin{proof}
Put \(s=t^2\).  Then
\[
  u=As+Bs^2+Cs^3+\BigO(s^4),
  \quad
  A=\kappa_2/2,\ B=\kappa_4/24,\ C=\kappa_6/720.
\]
Substitute
\[
  t=(u/A)^{1/2}(1+\alpha u+\beta u^2+\BigO(u^3)).
\]
Comparison in \(u=At^2+Bt^4+Ct^6+\cdots\) gives
\[
  \alpha=-\frac{B}{2A^2},
  \qquad
  \beta=\frac{7B^2-4AC}{8A^4}.
\]
Inserting \(A,B,C\) yields the stated coefficients.  The two signs are the
two branches supplied by \cref{thm:gq-centered-palindromic}.
\end{proof}

For \(P=G_{n,k}\), \cref{thm:gq-one-log} gives
\[
  \kappa_{2r}
  =\frac{B_{2r}}{2r}\Delta_{2r}(n,k),
  \qquad
  \kappa_2=\frac{d(n+1)}{12}.
\]
Thus the centered Gaussian inverse is an explicit
Bernoulli--Faulhaber Puiseux series.

\section{Cyclotomic roots and the alternating base}
\label{sec:gaussian-cyclotomic}

\begin{theorem}[Square-free Gaussian covering]
\label{thm:gq-square-free}
The cyclotomic factorization is
\[
  G_{n,k}(q)=\prod_{m=2}^{n}\Phi_m(q)^{e_m},
  \qquad
  e_m=\Floor{\frac nm}
      -\Floor{\frac km}
      -\Floor{\frac{n-k}{m}}.
\]
Every \(e_m\) is either zero or one.  Consequently every complex zero of
\(G_{n,k}\) is a simple root of unity, and the local inverse at a zero is an
ordinary Taylor branch rather than a Puiseux branch of higher order.
\end{theorem}

\begin{proof}
The factorization \(1-q^j=\prod_{m\mid j}\Phi_m(q)\) shows that the
multiplicity of \(\Phi_m\) in \(\QPochhammer{q}{q}{N}\) is
\(\Floor{N/m}\).
Subtracting the denominator multiplicities proves the formula for \(e_m\).
Write \(n=k+b\).  The elementary carry identity
\[
  \Floor{\frac{k+b}{m}}
  -\Floor{\frac{k}{m}}
  -\Floor{\frac{b}{m}}\in\{0,1\}
\]
proves square-freeness.  Distinct cyclotomic polynomials are coprime and each
is separable over characteristic zero.  Thus every zero is simple.  The
analytic inverse function theorem then supplies an ordinary local inverse.
\end{proof}

If \(\zeta\) is primitive of order \(m\) and \(e_m=1\), put
\[
  R_m(q):=\prod_{\substack{2\le r\le n\\r\ne m}}\Phi_r(q)^{e_r}.
\]
Then
\[
  G_{n,k}'(\zeta)=\Phi_m'(\zeta)R_m(\zeta)\ne0,
  \qquad
  q(y)=\zeta+
  \frac{y}{\Phi_m'(\zeta)R_m(\zeta)}+\BigO(y^2).
\]

\begin{proposition}[Complete first jet at \(q=-1\)]
\label{prop:gq-minus-one}
One has
\[
  G_{n,k}(-1)=
  \begin{cases}
  0,&n\text{ even and }k\text{ odd},\\[0.25em]
  \displaystyle
  \binom{\Floor{n/2}}{\Floor{k/2}},
  &\text{otherwise}.
  \end{cases}
\]
In every nonzero case, \(d=k(n-k)\) is even and
\[
  G_{n,k}'(-1)
  =-\frac d2\binom{\Floor{n/2}}{\Floor{k/2}}.
\]
In the zero case, writing \(n=2m\) and \(k=2r+1\),
\[
  G_{2m,2r+1}'(-1)=(m-r)\binom mr.
\]
Here \(r<m\), so this derivative is nonzero.  Thus \(-1\) has root
multiplicity one over \(\IntegerNumbers\), and after base change over every
characteristic-zero commutative ring.
\end{proposition}

\begin{proof}
Separate even and odd exponents in the product quotient and pair every
vanishing factor \(1-q^{2s}\) with a factor \(1-q^2\).  Since
\[
  \lim_{q\to-1}\frac{1-q^{2s}}{1-q^2}=s,
  \qquad
  \lim_{q\to-1}(1-q^{2s+1})=2,
\]
the surviving factorial ratios reduce to the displayed binomial coefficient.
When \(n\) is even and \(k\) odd, there is exactly one unpaired even numerator
factor, so the value is zero.

In every other parity case \(d=k(n-k)\) is even.  Differentiate the
reciprocity identity \(G(q)=q^dG(q^{-1})\) and set \(q=-1\).  This gives
\[
  G'(-1)=-dG(-1)-G'(-1),
\]
and hence the displayed even-degree derivative formula.

In that case the quotient after removing the unique factor \(1+q\) has limit
\[
  \frac{m!}{r!(m-r-1)!}=(m-r)\binom mr.
\]
Because \(1+q\) has derivative one at \(-1\), this limit is exactly
\(G_{2m,2r+1}'(-1)\).  The standing assumptions \(0<k<n\) imply \(r<m\),
so the integer slope is positive and the alternating zero is simple.  Root
multiplicity is preserved by the injective base change from
\(\IntegerNumbers\) to any characteristic-zero commutative ring.
\end{proof}

\begin{remark}[Exact Lean crosswalk]
The four parity values are proved over every commutative ring by
\lean{gaussianBinomial_neg_one_even_even},
\lean{gaussianBinomial_neg_one_odd_even},
\lean{gaussianBinomial_neg_one_odd_odd}, and
\lean{gaussianBinomial_neg_one_even_odd_eq_zero}.  The displayed exceptional
derivative is exactly
\lean{gaussianBinomial_derivative_eval_neg_one_even_odd}.  Its companion
\lean{gaussianBinomial_derivative_eval_neg_one_of_even_degree} gives the full
even-degree reciprocity derivative law.  Finally,
\lean{gaussianBinomial_even_odd_rootMultiplicity_int} and
\lean{gaussianBinomial_even_odd_rootMultiplicity} prove that the alternating
zero has multiplicity one when \(r<m\), respectively over \(\IntegerNumbers\)
and over every characteristic-zero commutative ring.  These declarations
belong to \path{GaussianBinomialAtNegOne.lean},
\path{QBinomialReciprocity.lean}, and
\path{GaussianBinomialAtNegOneDerivative.lean}; they do not assert the later
all-cyclotomic square-freeness theorem.
\end{remark}

\section{The reciprocal and large-target inverse}
\label{sec:gaussian-infinity}
\label{sec:gaussian-reciprocity}

Write the small-\(Q\) expansion
\[
  G_{n,k}(Q)=1+Q+c_2Q^2+c_3Q^3+\BigO(Q^4)
\]
with the \(c_2,c_3\) of \cref{sec:gaussian-zero}, and choose a branch
\(Y=y^{1/d}\).

\begin{theorem}[Correct large-target Gaussian inverse]
\label{thm:gq-large-inverse}
On every branch for which \(Y\to\infty\),
\begin{align*}
  q={}&Y-\frac1d
  +\left(\frac{d-1}{2d^2}-\frac{c_2}{d}\right)Y^{-1}\\
  &+\left(
    -\frac{(d-1)(d-2)}{3d^3}
    +\frac{(d-2)c_2}{d^2}
    -\frac{c_3}{d}
  \right)Y^{-2}
  +\BigO(Y^{-3}).
\end{align*}
\end{theorem}

\begin{proof}
Reciprocity gives
\[
  y=q^dG_{n,k}(q^{-1}).
\]
Set \(x=Y^{-1}\) and
\[
  q=Y(1+Ax+Bx^2+Cx^3+\BigO(x^4)).
\]
Then
\[
  q^{-1}
  =x-Ax^2+(A^2-B)x^3+\BigO(x^4),
\]
and therefore
\[
  G_{n,k}(q^{-1})
  =1+x+(c_2-A)x^2+
   (A^2-B-2Ac_2+c_3)x^3+\BigO(x^4).
\]
Divide the equation by \(Y^d\) and take logarithms.  Comparing the
coefficients of \(x,x^2,x^3\) in
\[
  0=d\log(1+Ax+Bx^2+Cx^3)
    +\log G_{n,k}(q^{-1})
\]
gives successively
\[
  A=-\frac1d,\qquad
  B=\frac{d-1}{2d^2}-\frac{c_2}{d},
\]
and
\[
  C=-\frac{(d-1)(d-2)}{3d^3}
    +\frac{(d-2)c_2}{d^2}
    -\frac{c_3}{d}.
\]
Multiplication by \(Y\) produces the displayed expansion and leaves a
remainder of order \(Y^{-3}\).
\end{proof}

\begin{remark}
Several source reports printed
\(-{(d-1)(2d-1)}/{(3d^3)}
+{2c_2(d-1)}/{d^2}-{c_3}/{d}\) at order \(Y^{-2}\).
The coefficient comparison above shows that expression is incorrect.  The
correct coefficient vanishes in the expected low-degree special cases and is
the one retained here.
\end{remark}

\section{Multinomial extension}
\label{sec:qmultinomial}

Throughout this section let \(s\ge1\), let
\(n_1,\ldots,n_s\in\IntegerNumbers_{\ge0}\), and put
\(n=n_1+\cdots+n_s>0\).  Set
\[
  M_{\bm n}(q)
  :=\QMultinomial{n}{n_1,\ldots,n_s}{q}
  =\frac{\QPochhammer{q}{q}{n}}
  {\QPochhammer{q}{q}{n_1}\cdots\QPochhammer{q}{q}{n_s}},
  \qquad
  D(\bm n):=\sum_{i<j}n_in_j,
\]
\[
  \Delta_m(\bm n)
  :=S_m(n)-\sum_{i=1}^{s}S_m(n_i).
\]
Thus $M_{\bm n}(q)$ is a section-local alias for the canonical
$q$-multinomial, while $D(\bm n)$ and $\Delta_m(\bm n)$ are section-local
auxiliaries only; none of the three symbols acquires a global meaning.

\begin{theorem}[Complete multinomial expansion at one]
\label{thm:qm-one-log}
For \(q=\EulerE^{-t}\) and \(|t|<2\pi/n\), take the holomorphic logarithm
of the normalized quotient whose value at \(t=0\) is zero.  Then
\[
  \log\frac{M_{\bm n}(\EulerE^{-t})}
                {\binom{n}{n_1,\ldots,n_s}}
  =-\frac{D(\bm n)}2t+
   \sum_{r\ge1}
   \frac{B_{2r}\Delta_{2r}(\bm n)}
        {2r(2r)!}\,t^{2r}.
\]
If \(D(\bm n)>0\), equivalently at least two of the \(n_i\) are positive,
then
\[
 \EulerE^{D(\bm n)t/2}
 \frac{M_{\bm n}(\EulerE^{-t})}{\binom{n}{n_1,\ldots,n_s}}
\]
is a nonconstant even function with a nondegenerate minimum at \(t=0\), and
its inverse has the two Puiseux branches of
\cref{prop:gq-centered-puiseux}.  If \(D(\bm n)=0\), then
\(M_{\bm n}\equiv1\) and no local inverse exists.
\end{theorem}

\begin{proof}
Apply \cref{lem:gq-bernoulli-factor} to the numerator
\(\QPochhammer{q}{q}{n}\) and every denominator
\(\QPochhammer{q}{q}{n_i}\).  Since \(|jt|<2\pi\) for every factor, the
normalized factor logarithms sum to the holomorphic branch specified in the
statement.  Constants give the ordinary multinomial
coefficient and the power sums give \(\Delta_{2r}(\bm n)\).  The linear
difference is
\[
  \frac12\left(n^2-\sum_i n_i^2\right)
  =\sum_{i<j}n_in_j=D(\bm n).
\]
It remains to justify the hypotheses for centered inversion.  Put
\(N_i=n_1+\cdots+n_i\).  The telescoping factorization
\[
 M_{\bm n}(q)=\prod_{i=2}^{s}\GaussianBinomial{N_i}{n_i}{q}
\]
proves polynomiality, with the empty product for \(s=1\) understood as one.
Boundary factors with \(n_i=0\) or \(N_{i-1}=0\) are one.  By
\cref{prop:gq-positive-palindromic}, every nontrivial Gaussian factor has
nonnegative coefficients, equal positive endpoints, and is palindromic;
hence so does their product.
Its degree is
\[
 \sum_{i=2}^{s}n_iN_{i-1}
 =\sum_{j<i}n_in_j=D(\bm n).
\]
Thus it is nonconstant exactly when \(D(\bm n)>0\), in which case
reciprocity makes the centered normalization even and
\cref{thm:gq-centered-palindromic,prop:gq-centered-puiseux} applies.  If
\(D(\bm n)=0\), at most one part is positive and the same factorization gives
\(M_{\bm n}=1\).
\end{proof}

\begin{theorem}[Cyclotomic carry and exact ramification]
\label{thm:qm-carry}
Under the standing hypotheses above, let \(m\ge1\), let \(\zeta\) be a
primitive \(m\)-th root of unity, and write
\(n_i=m\ell_i+r_i\) with \(0\le r_i<m\).  The multiplicity of \(\zeta\) as
a zero of \(M_{\bm n}\) is
\[
  e_m
  =\Floor{\frac nm}
   -\sum_{i=1}^{s}\Floor{\frac{n_i}{m}}
  =\Floor{\frac{r_1+\cdots+r_s}{m}}.
\]
Thus \(0\le e_m\le s-1\).  If \(e_m>0\), the inverse at target zero has
exactly \(e_m\) local Puiseux determinations, counted with their choices of
the \(e_m\)-th root.
\end{theorem}

\begin{proof}
As in \cref{thm:gq-square-free}, cyclotomic factorization of each factorial
gives the first expression.  Since
\[
  n=m\sum_i\ell_i+\sum_i r_i,
\]
substitution gives the residue-carry formula.  The residue sum lies between
zero and \(s(m-1)\), proving \(e_m\le s-1\).

Factor locally
\[
  M_{\bm n}(q)=(q-\zeta)^{e_m}U(q),
  \qquad U(\zeta)\ne0.
\]
After choosing an analytic \(e_m\)-th root of \(U\) in a sufficiently small
disk, the equation \(M_{\bm n}(q)=y\) is equivalent to
\[
  (q-\zeta)\,U(q)^{1/e_m}=\omega\,y^{1/e_m},
\]
where \(\omega\) runs through the \(e_m\)-th roots of unity.  The analytic
inverse function theorem applies to the left side, whose derivative at
\(\zeta\) is \(U(\zeta)^{1/e_m}\ne0\).  This produces precisely the stated
Puiseux determinations.
\end{proof}

For \(s=2\), the carry count is at most one, recovering the square-free
Gaussian case.  For many parts, the exact same cyclotomic mechanism produces
higher ramification.  Centered inversion at \(q=1\) and cyclotomic inversion
therefore represent genuinely different singular mechanisms even when they
occur in the same polynomial family.
